\chapter{Seed Instability and Where It Concentrates in the In-House Lab}
\label{ch:diagnosis}

The 30 REL\_GLOBAL seeds from which the teachers of
Section~\ref{sec:method-teachers} are drawn do not rank the same way on
the seen and the unseen conditions, and relative to their mean loss
they spread far more on the unseen ones. We measure how far they
diverge and locate where they do so, without isolating a cause.
Unless a sentence says otherwise, every number comes from one
development set of 135 executions and 1,427 natural rows that had been
scored many times before (Table~\ref{tab:seed-population}), so the
chapter reports development comparisons and post hoc analyses about
REL\_GLOBAL in the in-house lab.
% 근거: RECORD:126 (적용 범위); ORIG:PROPHET_REL_GLOBAL_SEED_POPULATION_20261118/
% FINAL_REPORT_KO.md:15 (135실행·1,427행, 집합별 실행·행). 이 30 seed가 4.4절 교사
% 후보와 같은 모집단이다(상위5 seed 50572, 91513, 71556, 68498, 37406, 같은 파일:33).

\section{How far the seeds diverge}
\label{sec:diag-size}

Table~\ref{tab:seed-population} summarizes the population. On
\texttt{A+C FLAT} the per-seed NLL had a mean of 0.142 and a seed
standard deviation of 0.086, and it ranged from 0.036 to 0.420; on
\texttt{A+C NESTED} it ranged from 0.023 to 0.404. For these 30
checkpoints the rankings were similar across the two layouts: a seed
that did well on one tended to do well on the other, with a Spearman
rank correlation of 0.886 between the FLAT and NESTED losses.
% 출처: ORIGINAL_TEXT/outputs/PROPHET_REL_GLOBAL_SEED_POPULATION_20261118/FINAL_REPORT_KO.md:37
% (ρ=.8861, FLAT–ORIGINAL_18 ρ=-.1310). atlas 압축본(inputs/drop/PROPHET_RESEARCH_TRIAL_ATLAS_
% 20261123.zip)을 sources/drop/unzipped에 푼 사본에서 대조했다.
Between \texttt{A+C FLAT} and
\texttt{ORIGINAL\_18} the correlation was $-0.131$. A trade between adaptation to the background and preservation of the
seen condition would make it clearly negative, so this value does not
support such a trade at the level of seeds. These correlations are
descriptive, without intervals.

The seeds are not alike on the seen condition either. Their standard
deviation on \texttt{ORIGINAL\_18} was 0.076, close to the 0.086 on
\texttt{A+C FLAT}, so in absolute terms the two conditions spread the
seeds about equally. Relative to the mean loss the picture differs: the
standard deviation was about 0.12 of the mean on \texttt{ORIGINAL\_18}
(0.076 against 0.631) and about 0.61 of it on \texttt{A+C FLAT} (0.086
against 0.142).
% 출처(ORIGINAL_18 SD .076415, C_ONLY .068273/.288810, B+C .052282/.193977, 최악 .763586):
%   atlas OT/PROPHET_MARKER_HEAD_INPUT_20261119/FINAL_REPORT_KO.md:19,25-27 (같은 30 checkpoint,
%   같은 135실행; FLAT/NESTED 값이 SEED_POPULATION:37과 일치) (diag.md 19a).
% 8회차 문장 P3-4: 10장 값(0.067/0.075, 다른 실행)을 앞당겨 쓰던 문장은 뺐다(10장 표에 있다).

Three selection scores are also listed, each fixed before the A+C
losses of these 30 seeds were read: S1, a seed's top-1 agreement with
the 30-model majority on the unlabeled A+C inputs, its own vote
included; S2, its NLL on \texttt{C\_ONLY}; and S3, its cross-entropy on
the last FIT minibatch of training. Each correlation compares the order
in which a score ranks the seeds with their order by A+C loss, lowest
loss first, so a positive value means that the seeds the score puts
first tend to have the lower loss.
% 근거: SEED_POPULATION FINAL_REPORT_KO.md:35 (FLAT 평균 .141949, SD .086001,
% 범위 .036339--.419831; NESTED 평균 .125527, SD .087071, 범위 .022585--.404273;
% FLAT--NESTED ρ=.8861; FLAT--ORIGINAL_18 ρ=-.1310). RECORD:62--64, 100.

\begin{table}[htbp]
\centering
\caption[The 30 REL\_GLOBAL seeds on the 135-run development set]{The 30 REL\_GLOBAL seeds on the 135-run development set. Top:
per-seed NLL, execution-weighted within each seed. Bottom: Spearman
correlation between the ranking of the seeds by a quantity and their
ranking by NLL on A+C, so that a positive value means the seeds ranked
first have lower A+C loss; CE is cross-entropy. S1--S3 are the
selection scores defined before the table. For the last three
conditions only the worst seed is recorded, not the best.}
\label{tab:seed-population}
\small
\begin{tabular}{@{}lrrrr@{}}
\toprule
Condition & Runs / rows & Mean NLL & Seed SD & Range over seeds \\
\midrule
\texttt{A+C FLAT}     & 31 / 341 & 0.142 & 0.086 & 0.036--0.420 \\
\texttt{A+C NESTED}   & 29 / 319 & 0.126 & 0.087 & 0.023--0.404 \\
\texttt{C\_ONLY}      & 27 / 297 & 0.103 & 0.068 & worst 0.289 \\
\texttt{B+C}          & 32 / 352 & 0.059 & 0.052 & worst 0.194 \\
\texttt{ORIGINAL\_18} & 16 / 118 & 0.631 & 0.076 & worst 0.764 \\
\midrule
Ranked by & & & with FLAT & with NESTED \\
\midrule
NLL on \texttt{A+C NESTED}   & & & 0.886    & --- \\
NLL on \texttt{ORIGINAL\_18} & & & $-0.131$ & --- \\
S1: agreement, unlabeled A+C & & & 0.877 & 0.915 \\
S2: NLL on \texttt{C\_ONLY}  & & & 0.115 & 0.151 \\
S3: last FIT minibatch CE    & & & 0.160 & 0.033 \\
\bottomrule
\end{tabular}
\end{table}
% 근거: SEED_POPULATION FINAL_REPORT_KO.md:15, 35, 37, 43 (B+C 평균 .059181,
% ORIGINAL_18 평균 .630865); C_ONLY 평균 .103086은
% ORIG:PROPHET_MARKER_TARGET_RULE_20261119/FOLLOWUP_REPORT_KO.md:22 (같은 30 seed, 같은 집합).
% SD·최악 seed 세 행: MARKER_HEAD_INPUT FINAL_REPORT_KO.md:25-27 (위 주석). 최솟값은 atlas에 없다.
% 부호 규약: 원 파일에 명시 없음. S1(합의도, 클수록 좋음)이 +.8768이고 보고서가 "강하게 연관"으로
%   읽으므로 순위 대 순위(양수 = 손실 낮은 seed를 앞에)로 추론했다(diag.md 19b).

Of the three scores, only S1 tracked A+C loss in this lab, at 0.877 on
FLAT and 0.915 on NESTED; neither S2 nor S3 exceeded 0.160 in either
layout (Table~\ref{tab:seed-population}).
% 출처: ORIGINAL_TEXT/outputs/PROPHET_REL_GLOBAL_SEED_POPULATION_20261118/FINAL_REPORT_KO.md:39
% (S1 .8768/.9150, S2 .1146/.1506, S3 .1604/.0331).
S1 and the losses it ranks come from the same development runs, so
the correlation is in-sample, and S1 is not a rule for choosing one
seed: the seed it ranked first failed preservation
(Section~\ref{sec:method-teachers}).
Section~\ref{sec:diag-agreement} returns to it.

S3 is not an NLL on held-out FIT runs, which was not computed for this
population. Neither the final training loss, nor labels from another
unseen condition, nor the loss on the seen condition recovered how the
seeds would rank on these background conditions
(Section~\ref{sec:related-underspec}).
% 원문: D'Amour et al. arXiv:2011.03395 v2 §2.1 "risk-equivalent near-optimal predictors" 가 출발점이라
%   seen 조건에서 동등하지 않으면 해당하지 않는다는 한정은 원문과 맞다. 확인됨.
% 근거: SEED_POPULATION FINAL_REPORT_KO.md:19--26 (S1--S4 정의, S3 = "전체 FIT 보류
% NLL이 아니며, 음성 대조"), :37 (S2 .1146/.1506, S3 .1604/.0331). RECORD:64.
% S4(고정 무작위 순열)는 무작위 대조라 표와 본문에서 뺐다.
% 출처: atlas OT/PROPHET_SEED_INPUT_AUDIT_20261118/FINAL_REPORT_KO.md:115 ("누출 없는 FIT 내부 보류
%   checkpoint/점수는 이 계보에서 찾지 못했다"; diag.md 19c). 앞선 5-seed 계보의 FIT 훈련 내 NLL 대
%   A+C NLL ρ=−0.80(n=5, :126)은 다른 seed 계보라 쓰지 않았다.

\section{Where the seeds disagree}
\label{sec:diag-agreement}

The places where the seeds part can be found without labels. On 1,029 of the
1,309 background rows all 30 seeds gave the same top-1 action, and
every one of these unanimous answers was correct. The rows on which the
seeds' top-1 predictions differed carried between 84.2\% and 91.2\%
of the NLL of the 30-model probability average, depending on the
condition, consistent with the use of disagreement among ensemble
members as a qualitative indicator of predictive
uncertainty~\cite{lakshminarayanan2016simple}.
% 원문: Lakshminarayanan et al. NIPS 2017 §3.4 "This disagreement provides another useful qualitative way
%   to evaluate predictive uncertainty" (각주 9: Σ_m KL(p_θm‖p_E)). "variation ... as a signal"을
%   원문대로 정성적 지표로 좁혔다.
Disagreement on unlabeled
environment inputs therefore marks, without labels, the places where
the seeds' computations part. The distillation of
Chapter~\ref{ch:method} is designed to reduce this disagreement on
unlabeled environment inputs, and Chapter~\ref{ch:results} tests it
under a sealed contract. Following the teachers on FIT inputs instead is
one of the controls (Section~\ref{sec:gen-methods}); on FIT views the
30 seeds gave the same top-1 action on 90.3\% of the 82,008 views,
which Section~\ref{sec:disc-rejected} takes up.

Unanimity is not a guarantee of correctness: on
\texttt{ORIGINAL\_18}, 8 of the 89 unanimous rows were wrong, and on
the background, subsets of five seeds were sometimes unanimous and
wrong. The S1 ranking orders single seeds by their agreement with the
majority, which is related to, but not the same as, the pairwise
agreement of Baek et al.~\cite{baek2022agreement}.
% 원문: Baek et al. NeurIPS 2022 §3.1 식 (1)--(2) Agr(h,h′) 는 두 모델 짝의 일치도, §4 식 (4) 는 짝별
%   ID·OOD 일치도의 probit 선형 회귀. 확인됨.
% 근거: ORIG:PROPHET_REL_GLOBAL_SEED_POPULATION_20261118/DISAGREEMENT_PREFLIGHT_20261119/
% REPORT_KO.md:5 (배경 1,029행 만장일치 모두 정답, 무작위 5개 만장일치는 오류 있음,
% ORIGINAL_18 89행 중 8행 오답), :9--15 (만장일치 247+240+243+299=1,029; 배경 1,309행),
% :25 (갈린 행의 30확률 평균 NLL 기여 FLAT 91.2%, NESTED 89.1%, C_ONLY 89.0%,
% B+C 84.2%), :33--37 (크기 5 조합). RECORD:101--102.
% S1 ρ .8768/.9150: SEED_POPULATION FINAL_REPORT_KO.md:37.
% FIT view 90.3% = 74,074/82,008 (개요 B표 일치). 82,008 = 4,824 자연 경계 +
% 77,184 전달 view (GENERIC_HEAD_CAPACITY_MATCHED FINAL_REPORT_KO.md:9, 11).
% FIT view는 3.1절의 natural·delayed view 정의를 따른다.

A separate diagnosis, which CD does not use, located most of the A+C
errors in one file decision. Every top-1 error of the 30 seeds on the
A+C rows chose the synthetic file copy (COPY, action 16 in
Appendix~\ref{app:actions}) where the label was the synthetic file
search (DISCOVERY, action 15), or the reverse; a
rule without parameters that reads the observed result for the current
target described that decision on 1,290 of 1,309 background rows; and
a head given the rule's inputs lowered the A+C loss on new runs but did
not preserve the seen condition and was not adopted.
Appendix~\ref{app:diagnosis} gives the analysis.
% 2026-11-26 지도교수 지시(본문 40~50쪽): VII-3~VII-5 와 표 VII.2 를 부록 F 로 옮기고 여기에는 요약만 둔다.

